\section{A smaller height exponent and an anchor of exponent 2364}

This argument starts from the frozen seventh-round source with builder
SHA256 \texttt{78b54bcde24ff8a91091e102a4734c49878699f071d9230e5ce7a70ffd2114d0}.
The upstream manuscript remains pinned at
\texttt{036cebe4f1a4aa7ca89820a9b48f3629b4c364d6}.
The new assertions below are explicit applications of the same inherited
geometric, modularity, isogeny-height, group, quaternion, pole, prime-pattern
and Lucas inputs. They do not independently establish those inputs or
reclassify any historical omission as a new gap.

The selected arithmetic change in the source is only
$S_0=3^{2790}\rightsquigarrow3^{2364}$. The physical anchor's packing
coefficient stays equal to three. The local conclusions consumed here
remain scalar order $28$, Lucas-label order $45$ and input-label order $12$.
A strengthened primary detector test, if selected elsewhere in this
proposal, has its own local/completeness proof and does not weaken any of
these conclusions. The current integrated source and this proof are bound
by the accompanying independent checker report.

\subsection{The precise inherited height inputs}

Put
\[
 \mathcal S=\{2,3,5,7,13,17\},\qquad B_{\mathcal S}=46410,
 \qquad C_{\rm cond}=2^{32}3^9.
\]
For $z\in\mathbb Q\setminus\{0,1,2401,-343/32\}$, let
$M_{\mathcal S}(z)$ be the squarefree product of primes outside
$\mathcal S$ having odd contact with one of
$\infty,0,1,2401,-343/32$, and set
$R=B_{\mathcal S}M_{\mathcal S}(z)$.
Thus $R\ge1$ is a squarefree rational integer. This $R$ is the radical
used in the global height proof; it is not a detector pole input or an
elliptic-curve index.

The finite-support reapplication in
\texttt{proofs/prior/round4/arithmetic\_round4.tex} and the conductor
refinement in \texttt{proofs/prior/round5/arithmetic\_round5.tex} retain
the source's geometric construction and (LR1)--(LR5). In particular, in
the non-CM case they supply the same surface $A$, field $L'$, primitive
newform of level $N$, and finite nebentypus of order $t$, with
\[
 d'=[L':\mathbb Q]\le2304R,\qquad
 N\le C_{\rm cond}R^2,\qquad \operatorname{rad}(N)\mid R,
 \qquad t\le2d'\le4608R.
\]
The last estimate is a character-order bound, not a coefficient-field
degree bound. The proofs of normality, the scalar-character corrections,
the companion at three used to bound the conductor at five, and the
conductor's primitive-level convention remain inherited and unchanged.

The actual source is
\texttt{manuscripts/quartic\_v10.tex},
\texttt{sec:h5-level-repair} and \texttt{thm:quartic-height19},
especially lines 8915--9073. Its subgroup $H=\ker\epsilon$ contains
$-I$ and gives a modular curve $X_H$ over $\mathbb Q$, with
\[
 d_H=[\operatorname{SL}_2(\mathbb Z):\Gamma_H]
     =Nt\prod_{p\mid N}(1+1/p),\qquad
 \deg_B(X_H)\le d_H,\qquad g\le1+d_H/12.
\]
Here $g=\dim\operatorname{Jac}(X_H)$. The inherited realization,
absolute simplicity and Faltings--Poincare argument put the
two-dimensional surface $A$ in that Jacobian up to isogeny over $L'$;
therefore $g\ge2$. No bounded coefficient-field degree, squarefree level,
or trivial nebentypus is being introduced.

\subsection{An elementary Euler-product refinement}

Let
\[
 \mathcal P=\{2,3,5,7,11,13,17,19,23,29,31,37,41,43,47\}.
\]
These are exactly the primes below $53$. For
$a_p=(1+1/p)p^{-1/192}$, the following finite integer inequalities hold:
\begin{align*}
 &(p+1)^{192}>p^{193} &&(p\in\mathcal P),\\
 &54^{192}<53^{193},\\
 &\left(9\prod_{p\in\mathcal P}(p+1)\right)^{192}
       <32^{192}\left(\prod_{p\in\mathcal P}p\right)^{193}.
\end{align*}
The companion certificate verifies every integer comparison, as well as
the completeness and primality of the displayed list. The function
$(1+1/x)^{192}/x$ is decreasing for $x>0$, as a product of positive
decreasing functions. Hence $a_p\le1$ for every prime $p\ge53$,
whereas every displayed small-prime factor exceeds one. For every
positive integer $N$ it follows that
\[
 \prod_{p\mid N}(1+1/p)
 =\operatorname{rad}(N)^{1/192}\prod_{p\mid N}a_p
 \le\operatorname{rad}(N)^{1/192}\prod_{p\in\mathcal P}a_p
 <\frac{32}{9}\operatorname{rad}(N)^{1/192}.
\]
This controls arbitrary subsets of the small primes; a product of
unconditionally smaller-than-one factors has not been inserted.

Define
\[
 \alpha=18+\frac1{32}=\frac{577}{32},\qquad
 T=2^{14}C_{\rm cond}R^{577/192}.
\]
The level, character and radical bounds above now give
\[
 d_H<4608\,\frac{32}{9}C_{\rm cond}R^{3+1/192}=T.
\]
The number $53$ is only a cutoff in a numerical Euler-product estimate.
It does not enlarge $\mathcal S$, discard contacts at smaller primes,
or alter any local logarithm, pole, ramification or prime-selection
hypothesis.

\subsection{Reapply the Jacobian and isogeny-height bounds}

The genus estimate can be used without the source's coarse replacement
$g\le d_H$. Indeed $g\ge2$ and $g\le1+d_H/12$ imply $d_H\ge12$,
and hence
\[
 g\le1+d_H/12\le d_H/6<T/6.
\]
Use the same stable Faltings metric as the source. Its cited
Javanpeykar bound is
\[
 h_F(J_H)\le13000000\,\deg_B(X_H)^5g
       <\frac{13000000}{6}T^6.
\]
The source's Gaudron--R\'emond application to the principally polarized
$J_H$ and its isogeny factor $A\times Z$ gives
\[
 h_F(A)+h_F(Z)\le h_F(J_H)+V,
\]
\[
 V=2g^2\left(8g^2\log(7g)+\log d'
        +\log\max(1,\log d',h_F(J_H))\right).
\]
Retain its already proved upper estimate, with $u=\log T$:
\[
 V\le16T^4(2+u)+2T^2(25+7u)\le T^6.
\]
All hypotheses used for this estimate still hold: $T\ge16$,
$g<T$, $\log d'<8+\log R\le8+u$, and the new Jacobian bound is
smaller than the source's $2^{24}T^6$ bound. The elementary inequalities
$2+\log T\le T^2/32$ and $25+7\log T\le T^4/4$ hold at $16$
and persist by differentiation, exactly as in the source. No new
height-variation theorem is used.

The retained Bost lower bound gives $h_F(Z)>-3g/2>-3T/2$. Since
$3T/2<T^6$ for $T\ge16$ and
\[
 \frac{13000000}{6}+2<2^{22},
\]
we obtain
\[
 h_F(A)<2^{22}T^6
       =2^{106}C_{\rm cond}^6R^{\alpha}.
\]
The original CM estimate $h_F(A)<854R$ is smaller than this bound
as well. Applying the same geometric comparison
$h(z)\le2\max(0,h_F(A))+128$ gives
\[
 h(z)<2^{107}C_{\rm cond}^6R^{\alpha}+128.
\]
Finally the exact comparison
$2^{32}B_{\mathcal S}<3^{32}$ implies
$B_{\mathcal S}^{1/32}<3/2$. Thus, with the integer
\begin{equation}\label{r8:height-constant}
 \mathcal H=3\cdot2^{106}C_{\rm cond}^6B_{\mathcal S}^{18}+128,
\end{equation}
and using $M_{\mathcal S}(z)\ge1$, the reapplication proves
\begin{equation}\label{r8:height}
 \boxed{h(z)\le\mathcal H\,M_{\mathcal S}(z)^{577/32}.}
\end{equation}
The harmless weak inequality is convenient for all subsequent consumers.
The exact integer $\mathcal H$ has bit length $665$ and satisfies
$\mathcal H<3\cdot2^{663}$. If $H_7$ denotes the preceding
683-bit constant, then
\[
 H_7-128=247520(\mathcal H-128).
\]
This smaller exponent and coefficient follow from the displayed
inequalities in the inherited proof, rather than from a formal
substitution into its final exponent-nineteen theorem statement.

\subsection{The retained even anchor and its noncircular local order}

Write $K=\mathbb Q(\rho)$, $\rho^2=2$ and
$N_K(x+y\rho)=x^2-2y^2$. Let the four root labels be $(n,a,t,b)$,
and put
\[
 C=\prod_i x(Q_i)=C_0+C_1\rho,\qquad B_0=4096C_0.
\]
The scalar point is $S=[a]P$, its shifted point is
$Q=S-P=[a-1]P$, and the retained chart and nonzero guard give
$a\ne0$, $a-1\ne0$, and an existing field coordinate
$w=h_S\tau$ satisfying $h_S^2=x(S)$ and $\tau^2=f(Q)$.
Set
\[
 S_0=3^{2364},\qquad
 \Theta=N_K(w)+3S_0N_K(B_0+X\rho),
\]
and retain the three-witness conic
\[
 u^2+6v^2-2z^2-3\Theta^2+2S_0^2=0.
\]
Its exact point/input domain is
\begin{equation}\label{r8:anchor-domain}
 \text{all four labels odd},\qquad
 X\in\mathbb Z_{(2)}\cap\mathbb Z_{(3)},\qquad
 0\ne a-1\in2\cdot3^{1181}\mathbb Z.
\end{equation}

Here is the local order, making explicit why the parity argument is
not circular. The same quaternion criterion, including its zero
argument, tests $2$- and $3$-integrality of
\[
 \Theta/S_0=N_K(w/3^{1182})+3N_K(B_0+X\rho).
\]
At $3$ the norm is anisotropic modulo $3$, so the two nonzero summands
have opposite valuation parity and must separately be integral.
This gives integrality of $B_0,X$ and depth at least $1182$ for both
coefficients of $w$. Before assuming oddness, the retained leading-residue
proof gives
\[
 v_3(C_0)=-\sum_i e(n_i),\qquad
 e(n)=0\ (n\text{ odd}),\quad e(n)=2+2v_3(n)\ (n\text{ even}).
\]
Consequently all labels are odd. Only then use that $h_S$ is a
$3$-adic unit and $v_3(w)=1+v_3(a-1)$ to infer the displayed
$3$-divisibility.

At $2$, with $v(2)=1$ and $v(\rho)=1/2$, the retained integral formal
series gives $v(z([j]P))=3/2+v_2(j)$ for nonzero integers $j$.
The already established oddness implies
$v(h_S)=-3/2$ and $v(w)=v_2(a-1)\ge1$, so both rational
coefficients of $w$ are $2$-integral. Also $B_0$ is $2$-integral.
The only possibly nonintegral term left in $\Theta/S_0$ is $-6X^2$;
its negative valuation cannot cancel the integral terms. This gives
$X\in\mathbb Z_{(2)}$. The same valuations and the quaternion
criterion prove the converse. Zero norm summands and all chart sign
choices are covered, and no new square-root branch is introduced.

The supplying arguments are the packed-anchor leading-residue proof,
\texttt{proofs/prior/round6/arithmetic\_round6.tex}, and the actual
upstream formal/quaternion passages at lines 8132--8244 and
10973--11024. The conic witnesses are existentially equivalent to
\eqref{r8:anchor-domain}; no bijection between all witness tuples is
claimed.

\subsection{Apply the new height bound to the actual consumers}

The current input and three scalar contacts are
\[
 \beta=2401+2754(2X+1)^2
       =11016X^2+11016X+5155,
 \qquad r=\beta(\beta-1)(32\beta+343).
\]
The anchor excludes $X=-1/2$, so $\beta$ lies in the domain of
\eqref{r8:height}. The current scalar proof supplies
$M_{\mathcal S}(\beta)^{28}\mid a-1$. With $c=|a-1|$,
\[
 h(\beta)\le\mathcal H c^{577/896}.
\]
The detector poles retain orders $1,12,45$, their current contact-unit
normalization, and the same three positive square scales
$1,C_0^2,(C_0+C_1)^2$. Their unchanged fractional-linear map
$\phi(t)=-49+1/t$ is defined and avoids the four finite target contacts
on the proved source domain. The long logarithm $G_{12}$ still supports
order $45$ and $G_2$ supports order $12$ at all consumed primes.
These are the actual local arguments of the preceding proof and do not
depend on the global exponent being nineteen.

For completeness, if the three contact radicals divide a nonzero integer
$N$ to order $k$, their individual height estimates are now bounded by
$T_0=(\mathcal H+\log50)|N|^{\alpha/k}$.
The retained scale argument gives
$h(C)\le3T_0+(5/2)\log2<7\mathcal H|N|^{\alpha/k}$,
because $3\log50+(5/2)\log2<16$ and $\mathcal H\ge4$.
All four labels are already odd, so the inherited product-height
identity and root lower bound give
\begin{equation}\label{r8:transfer}
 B^2-6<7\mathcal H|N|^{\alpha/k},\qquad
 B=\max(|n|,|a|,|t|,|b|).
\end{equation}
This is the product in $K$, not its rational norm. The optional
projective-height improvement is not needed for this selection.

\subsection{Exact numerical margins, with a weaker sufficient scalar target}

Put
\[
 c_*=2\cdot3^{1181},\qquad b_*=c_*-1.
\]
Every solution of the selected anchor has $c\ge c_*$ and
$|a|\ge b_*$. We use $h(\beta)<c/4$, in place of the former
stronger $c/80$ target, and check the actual final input contradiction
below. The exact inequalities
\[
 3^{53}>2^{84},\qquad \mathcal H<3\cdot2^{663}
\]
give a short transparent certificate. If $L=\log_2 3$, then
$\log_2c\ge1+1181L$ and
$\log_2(4\mathcal H)<665+L$. The difference is positive because
\[
 319+(319\cdot1181-896)\frac{84}{53}-665\cdot896
   =\frac{8199}{53}>0.
\]
Hence
\begin{equation}\label{r8:scalar-margin}
 c^{319/896}>4\mathcal H,\qquad
 h(\beta)<c/4\le(|a|+1)/4.
\end{equation}
The independent checker also verifies directly
$c_*^{319}>(4\mathcal H)^{896}$ using integers only.

For the literal Lucas polynomial
\[
 F=t^2-a(a-2)tb+(a-2)^2b^2-1,
\]
the retained bound is $|F|<5B^4$ for $B>4$. If $F\ne0$, apply
\eqref{r8:transfer} at order $45$. Since $\alpha/45<1/2$,
\[
 B^{143/360}<7\mathcal H5^{\alpha/45}+6<24\mathcal H.
\]
But $B\ge|a|\ge c/2\ge3^{1181}$ and
$24\mathcal H<9\cdot2^{666}$ give the opposite strict inequality:
\[
 (143\cdot1181-720)\frac{84}{53}-666\cdot360
       =\frac{1418412}{53}>0.
\]
Equivalently the checker verifies
$b_*^{143}>(24\mathcal H)^{360}$. Thus $F=0$.
The nonidentity root for $b$ gives $b\ne0$, so the inherited Lucas
theorem now gives $\log B>|a|-3$.

If $n\ne\beta$, its nonzero reduced numerator $N$ satisfies
$|N|\le2e^{h(\beta)}B$, and the local input comparison gives order
$12$. Applying \eqref{r8:transfer} and using $\alpha/12<2$ yields
\[
 B^{2-\alpha/12}<40\mathcal H e^{(\alpha/12)h(\beta)},
\]
\[
 \log B<\frac{12}{24-\alpha}\log(40\mathcal H)
                 +\frac{\alpha}{24-\alpha}h(\beta).
\]
The first inequality follows, for example, from
$7\mathcal H2^{\alpha/12}+6<40\mathcal H$; the added constant
is absorbed since $h(\beta)\ge0$ and $B>1$.
The same exact height bound gives $40\mathcal H<2^{670}$.
Since $100\cdot670<b_*$, we have
$\log(40\mathcal H)<|a|/100$. Finally $\alpha<19$ gives
\[
 \log B<\frac{12}{5}\frac{|a|}{100}
          +\frac{19}{5}\frac{|a|+1}{4}
       =\frac{487}{500}|a|+\frac{19}{20}<|a|-3,
\]
where the last inequality holds for $|a|>152$. The anchor supplies
that bound. This contradicts the Lucas conclusion. Every use of the
old $c/80$ margin in the final soundness proof has therefore been
replaced by a checked use of \eqref{r8:scalar-margin}; the old numerical
inequality is not asserted for this smaller anchor.

Consequently $n=\beta\in\mathbb Z$. Together with the anchor's
$2$- and $3$-integrality, this forces $X\in\mathbb Z$:
$11016=2^3 3^4 17$, and a negative valuation of $X$ at any other
prime makes the quadratic term of $\beta$ uniquely minimal, including
the leading-coefficient valuation one at $17$.

\subsection{Completeness, implementation, and verification boundary}

The new height estimate is used only in soundness. It changes no point,
scale, pole, logarithm, detector contact, or congruence in completeness.
The even anchor requires only the finite preparation
$3^{1181}\mid a-1$, while the inherited construction allows arbitrary
divisibility of the positive scalar index $a-1$. Retain its order of
choices: scalar preparation and roots, fixed positive square scales,
continuity conditions, then separated contact primes and the primary
progression $\Pi^{44}\mid h$. The scalar pole preparation, Gaussian
positivity witnesses, both-place formal arguments and local avoidance
conditions remain those of the selected local proofs.

No coordinate is added by changing the fixed integer $S_0$.
Its even packing coefficient remains three, so the symbolic physical
anchor equations are unchanged. The convenient identity
$3^{2364}=(3^{394})^6$ supports a separately audited input-fixed
diagonal coefficient change. It is that separate algebraic certificate,
not this height argument, which proves any reported coefficient maximum.

The finite checker verifies the source anchor and input equations, local
seed calculations, Euler-product certificate, explicit constants, and
all three selected numerical consumers. The complete scalar/detector
proofs and the inherited global arithmetic still require their stated
mathematical review. This is a conditional reapplication with a smaller
height exponent and an explicit valid parameter choice; no optimality,
whole-proof formalization, or historical priority claim is made.
